% Part: proof-theory
% Chapter: sequent-calculus
% Section: proof-examples

\documentclass[../../../include/open-logic-section]{subfiles}

\begin{document}

\olfileid[pt]{pt}{seq}{ex}

\olsection{Exemplos de demonstrações}

Para construir !!{proof}s de sequentes, em geral é conveniente trabalhar
de baixo para cima a partir do sequente final, escolhendo nele !!a{formula}
como !!{formula} ativa e escrevendo as premissas da regra correspondente
acima do sequente; repete-se então o processo até chegar a axiomas. Por
exemplo, suponha que queiramos demonstrar o sequente
\[(!C \land !D) \lif !E \Sequent !C \lif (!D \lif !E)\]
Considere $!C \lif (!D \lif !E)$: ela tem a forma $!A \lif !B$ e
ocorre à direita do sequente. Fazendo o sequente inteiro corresponder
à regra $\RightR{\lif}$ de \Log{G1c},
\[\Axiom$ !A, \Gamma \fCenter \Delta, !B$
\RightLabel{\RightR{\lif}}
\UnaryInf$ \Gamma \fCenter \Delta, !A \lif !B $
\DisplayProof\]
somos levados a tomar $\Gamma = \{!C \lif (!D \lif !E)\}$, $\Delta =
\emptyset$, $!A \ident !C$ e $!B \ident !D \lif !E$. Assim, a última
inferência de uma demonstração poderia ser
\[\Axiom$ !C, (!C \land !D) \lif !E \fCenter !D \lif !E$
\RightLabel{\RightR{\lif}}
\UnaryInf$(!C \land !D) \lif !E \fCenter !C \lif (!D \lif !E)$
\DisplayProof\]
Se repetirmos essa ideia, tomando $!D \lif !E$ como fórmula ativa,
obteremos
\[\Axiom$ !D, !C, (!C \land !D) \lif !E \fCenter !E$
\RightLabel{\RightR{\lif}}
\UnaryInf$!C, (!C \land !D) \lif !E \fCenter !D \lif !E$
\RightLabel{\RightR{\lif}}
\UnaryInf$(!C \land !D) \lif !E \fCenter !C \lif (!D \lif !E)$
\DisplayProof\]
Agora precisamos nos concentrar em $(!C \land !D) \lif !E$ e usar a regra $\LeftR{\lif}$ de \Log{G1c},
\[\Axiom$ \Gamma \fCenter \Delta, !A$
\Axiom$ !B, \Gamma \fCenter \Delta$
\RightLabel{\LeftR{\lif}}
\BinaryInf$ !A \lif !B, \Gamma \fCenter \Delta$
\DisplayProof\]
Isso nos leva a
\[\Axiom$!D, !C \fCenter !E, !C \land !D$
\Axiom$!D, !C, !E \fCenter !E$
\RightLabel{\LeftR{\lif}}
\BinaryInf$ !D, !C, (!C \land !D) \lif !E \fCenter !E$
\RightLabel{\RightR{\lif}}
\UnaryInf$!C, (!C \land !D) \lif !E \fCenter !D \lif !E$
\RightLabel{\RightR{\lif}}
\UnaryInf$(!C \land !D) \lif !E \fCenter !C \lif (!D \lif !E)$
\DisplayProof\]
Usando a !!{formula} $!C \land !D$ no sequente superior esquerdo como
principal e a regra $\RightR{\land}$, obtemos
\[
\Axiom$!D, !C \fCenter !E, !C$
\Axiom$!D, !C \fCenter !E, !D$
\RightLabel{\RightR{\land}}
\BinaryInf$!D, !C \fCenter !E, !C \land !D$
\Axiom$!D, !C, !E \fCenter !E$
\RightLabel{\RightR{\lif}}
\BinaryInf$ !D, !C, (!C \land !D) \lif !E \fCenter !E$
\RightLabel{\RightR{\lif}}
\UnaryInf$!C, (!C \land !D) \lif !E \fCenter !D \lif !E$
\RightLabel{\RightR{\lif}}
\UnaryInf$(!C \land !D) \lif !E \fCenter !C \lif (!D \lif !E)$
\DisplayProof
\]
Até aqui, tudo bem. Observe que tudo o que fizemos também funciona em \Log{G3c},
pois as regras usadas até agora são as mesmas em \Log{G1c}
e~\Log{G3c}. Chegamos a um fragmento de !!a{proof} no qual todos os sequentes
superiores contêm a mesma !!{formula} à esquerda e à direita. Contudo,
eles ainda não são propriamente axiomas.

Em \Log{G1c}, precisamos construir !!{proof}s dos sequentes superiores a partir
de axiomas da forma $!A \Sequent !A$. Isso é feito facilmente com as
regras de enfraquecimento; por exemplo, o sequente superior mais à esquerda $!D, !C \Sequent
!E, !C$ pode ser !!{prove}d
da seguinte maneira:
\[
\Axiom$!C \fCenter !C$
\RightLabel{\LeftR{\Weakening}}
\UnaryInf$!D, !C \fCenter !C$
\RightLabel{\RightR{\Weakening}}
\UnaryInf$!D, !C \fCenter !E, !C$
\]
No conjunto, temos !!a{proof} em \Log{G1c} com o seguinte aspecto:
\[
\Axiom$!C \fCenter !C$
\RightLabel{\LeftR{\Weakening}}
\UnaryInf$!D, !C \fCenter !C$
\RightLabel{\RightR{\Weakening}}
\UnaryInf$!D, !C \fCenter !E, !C$
\Axiom$!D \fCenter !D$
\RightLabel{\LeftR{\Weakening}}
\UnaryInf$!D, !C \fCenter !D$
\RightLabel{\RightR{\Weakening}}
\UnaryInf$!D, !C \fCenter !E, !D$
\RightLabel{\RightR{\land}}
\BinaryInf$!D, !C \fCenter !E, !C \land !D$
\Axiom$!E \fCenter !E$
\RightLabel{\LeftR{\Weakening}}
\UnaryInf$!C, !E \fCenter !C$
\RightLabel{\LeftR{\Weakening}}
\UnaryInf$!D, !C, !E \fCenter !E$
\RightLabel{\LeftR{\lif}}
\BinaryInf$ !D, !C, (!C \land !D) \lif !E \fCenter !E$
\RightLabel{\RightR{\lif}}
\UnaryInf$!C, (!C \land !D) \lif !E \fCenter !D \lif !E$
\RightLabel{\RightR{\lif}}
\UnaryInf$(!C \land !D) \lif !E \fCenter !C \lif (!D \lif !E)$
\DisplayProof
\]
A estrutura dessa demonstração (em particular, sua altura) independe
de quais sejam as !!{formula}s $!C$, $!D$ e~$!E$. Podemos substituir
esses símbolos na !!{proof} acima por quaisquer !!{formula}s, e o
  resultado ainda será !!{proof} em~\Log{G1c}. A demonstração em \Log{G1c}
é \emph{esquemática}.

Em \Log{G3c}, os axiomas são sequentes da forma $!A, \Gamma \Sequent
\Delta, !A$, em que $!A$ deve ser atômica. Assim, embora $!D, !C \Sequent
!E, !C$ pareça um axioma de \Log{G3c} ($\Gamma = \{!D\}$; $\Delta =
\{!E\}$), ele só é \emph{realmente} um axioma se $!C$ for de fato atômica.
Nosso fragmento de demonstração só é uma demonstração completa em \Log{G3c} se $!C$,
$!D$ e $!E$ forem !!{formula}s atômicas.

Embora, a rigor, não tenhamos mostrado que
\[\Log{G1c} \Proves (!C \land !D) \lif !E \fCenter !C \lif (!D \lif
!E),\] na prática isso é tudo o que precisamos (e podemos) fazer. Isso
ocorre porque todo sequente da forma $!A, \Gamma \Sequent \Delta, !A$ é
demonstrável em \Log{G3c}, mesmo que $!A$ não seja atômica. Podemos mostrar
isso por indução sobre a profundidade~$\depth{!A}$ de~$!A$.

\begin{defn}\ollabel{defn:depth}
A \emph{profundidade}~$\depth{!A}$ de !!a{formula} é definida da seguinte maneira:
\begin{enumerate}
  \item Se $!A$ é atômica, $\depth{!A} = 0$.
  \item Se $!A \ident \lnot !B$, $!A \ident \lforall[x][A]$ ou 
  $!A \ident \lexists[x][A]$, então $\depth{!A} = \depth{!B} + 1$.
  \item Se $!A \ident (!B \land !C)$, $(!B \lor !C)$ ou $(!B \lif
  !C)$, então $\depth{!A} = \max(\depth{!B},\depth{!C}) + 1$.
\end{enumerate}
\end{defn}

Não é difícil provar que $\depth{A(x)} = \depth{A(t)}$ para qualquer
termo~$t$. O fato importante para nós é que, em qualquer regra, a profundidade
da !!{formula} principal é sempre maior que a profundidade das
!!{formula}s ativas. Por exemplo, temos:
\begin{align*}
\depth{P(x)} = \depth{P(a)} & = 0,\\
\depth{\lforall[x][P(x)]} & = 1,\\
\depth{(\lforall[x][P(x)] \lif P(a))} & = \max(1,0) + 1 = 2.
\end{align*}
Podemos usar isso para provar alguns resultados por indução sobre~$\depth{!A}$,
como o seguinte.

\begin{prop}\ollabel{prop:G1c-axioms}
Para toda~$!A$, existe !!a{proof} em $\Log{G1c}$ de $!A \Sequent !A$
na qual todos os axiomas são atômicos.
\end{prop}

\begin{proof}
  Por indução sobre~$\depth{!A}$. Se $\depth{!A} = 0$, então $!A$ é
  atômica, e $!A \Sequent !A$ já é uma !!{proof} em~\Log{G1c} da
  forma exigida. Se $\depth{!A} > 0$, $!A$ é construída a partir de
  !!{formula}s de menor profundidade, por exemplo, $!A \ident \lnot !B$, $!A \ident
  (!B \land !C)$ ou $!A \ident \lforall[x][A(x)]$. A hipótese de
  indução é que, para toda $!D$ com $\depth{!D} < \depth{!A}$,
  $\Log{G1c} \Proves !D \Sequent !D$ a partir de axiomas atômicos.

  Se $!A \ident \lnot !B$, então $\depth{!B} < \depth{!A}$ e, pela
  hipótese de indução, existe !!a{proof}~$\pi_1$ em \Log{G1c}, a partir de
  axiomas atômicos, de
  $!B \Sequent !B$. Podemos estendê-la a !!a{proof}
  de $\lnot !B \Sequent \lnot !B$ da seguinte maneira:
  \[
  \AxiomC{}
  \RightLabel{$\pi_1$}
  \Deduce$!B \fCenter !B$
  \RightLabel{\LeftR{\lnot}}
  \UnaryInf$\lnot !B, !B \fCenter $
  \RightLabel{\RightR{\lnot}}
  \UnaryInf$\lnot !B \fCenter \lnot !B$
  \DisplayProof
\]

  Se $!A \ident (!B \land !C)$, então $\depth{!B} < \depth{!A}$ e
  $\depth{!C} < \depth{!A}$. Pela hipótese de indução, existem
  !!{proof}s~$\pi_1$ e~$\pi_2$ em \Log{G1c}, a partir de axiomas atômicos, de $!B
  \Sequent !B$ e $!C \Sequent !C$. Podemos estendê-las
  a !!a{proof} de $!B \land !C \Sequent !B \land !C$
  da seguinte maneira:
  \[
  \AxiomC{}
  \RightLabel{$\pi_1$}
  \Deduce$!B, \fCenter !B$
  \RightLabel{\LeftR{\Weakening}}
  \UnaryInf$!B, !C, \fCenter !B$
  \RightLabel{\LeftR{\land}}
  \UnaryInf$!B \land !C, !C\fCenter !B$
  \RightLabel{\LeftR{\land}}
  \UnaryInf$!B \land !C, !B \land !C \fCenter !B$
  \AxiomC{}
  \RightLabel{$\pi_2$}
  \Deduce$!C, \fCenter !C$
  \RightLabel{\LeftR{\Weakening}}
  \UnaryInf$!B, !C, \fCenter !C$
  \RightLabel{\LeftR{\land}}
  \UnaryInf$!B \land !C, !C \fCenter !C$
  \RightLabel{\LeftR{\land}}
  \UnaryInf$!B \land !C, !B \land !C \fCenter !C$
  \RightLabel{\RightR{\land}}
  \BinaryInf$!B \land !C \fCenter !B \land !C$
  \RightLabel{\LeftR{\Contraction}}
  \UnaryInf$!B \land !C \fCenter !B$
  \DisplayProof
  \]

  Agora, suponha $!A \ident \lforall[x][B(x)]$. Tome !!a{constant}~$c$
  que não ocorra em $\lforall[x][B(x)]$. Então $\depth{B(c)} =
  \depth{!B(x)} < \depth{!A}$ e, pela hipótese de indução, existe
  !!a{proof}~$\pi_1$ em \Log{G1c}, a partir de axiomas atômicos, de $!B(c)
  \Sequent !B(c)$. Podemos estendê-la a !!a{proof} de
  $\lforall[x][B(x)] \Sequent \lforall[x][B(x)]$ da
  seguinte maneira:
  \[
  \AxiomC{}
  \RightLabel{$\pi_1$}
  \Deduce$!B(c) \fCenter !B(c)$
  \RightLabel{\LeftR{\lforall}}
  \UnaryInf$\lforall[x][!B(x)] \fCenter !B(c)$
  \RightLabel{\RightR{\lforall}}
  \UnaryInf$\lforall[x][!B(x)] \fCenter \lforall[x][!B(x)]$
  \DisplayProof
\]

  Os casos restantes ficam como exercícios.
\end{proof}

\begin{prob}
  Complete a demonstração de \olref{prop:G1c-axioms}, tratando os
  casos em que $!A \ident (!B \lor !C)$, $!A \ident (!B \lif !C)$ e
  $!A \ident \lexists[x][B(x)]$.
\end{prob}

Para nossos fins, também poderíamos ter usado a !!{proof}
\[
\AxiomC{}
\RightLabel{$\pi_1$}
\Deduce$!B \fCenter !B$
\RightLabel{\RightR{\lnot}}
\UnaryInf$\fCenter !B, \lnot !B$
\RightLabel{\LeftR{\lnot}}
\UnaryInf$\lnot !B \fCenter \lnot !B$
\DisplayProof
\]
no caso em que $!A \ident \lnot !B$. Contudo, há sistemas de sequentes
nos quais isso não funcionaria. O cálculo de sequentes intuicionista
\Log{G1i} é semelhante a \Log{G1c}, exceto por restringir todos os
sequentes a, no máximo, uma !!{formula} no lado direito. Se
$\Delta = \emptyset$, então a primeira !!{proof} também seria
!!a{proof} em~\Log{G1i}, mas a segunda não seria, pois
um sequente contém tanto $!B$ quanto~$\lnot !B$ à direita.

Observe que, mesmo em \Log{G1c}, não poderíamos ter usado uma ordem alternativa
das regras no caso em que $!A \ident \lforall[x][!B(x)]$. O que segue
\emph{não} é !!a{proof}
  \[
  \AxiomC{}
  \RightLabel{$\pi_1$}
  \Deduce$!B(c) \fCenter !B(c)$
  \RightLabel{\RightR{\lforall}}
  \UnaryInf$!B(c) \fCenter \lforall[x][!B(x)]$
  \RightLabel{\LeftR{\lforall}}
  \UnaryInf$\lforall[x][!B(x)] \fCenter \lforall[x][!B(x)]$
  \DisplayProof
\]
pois a condição da variável própria em~$\RightR{\lforall}$ é violada.

\begin{prop}\ollabel{prop:G3c-axioms}
Todo sequente da forma $!A, \Gamma
\Sequent \Delta, !A$ (com $!A$ \emph{não} necessariamente atômica) é
demonstrável em~\Log{G3c}.
\end{prop}

\begin{proof}
A demonstração é muito semelhante à de \olref{prop:G1c-axioms}, mas
precisamos ter cuidado com a hipótese de indução. O caso em que $n =
\depth{!A} = 0$ é novamente fácil. Se $n=0$, $!A$ é atômica, e $!A,
\Gamma \Sequent \Delta, !A$ é um axioma de~\Log{G3c}.

Agora, suponha $n>0$. Novamente distinguimos casos. A hipótese de
indução é: para todo $D!$ com $\depth{!D} < n$, e todos $\Pi$
e $\Lambda$, $\Log{G3c} \Proves !D, \Pi \Sequent \Lambda, !D$. Eis
o caso de $!A \ident (!B \land !C)$. Usamos a hipótese de
indução duas vezes: uma para $!D = !B$, $\Pi = !C, \Gamma$ e $\Lambda
= \Delta$, obtendo !!a{proof}~$\pi_1$ de $!B, !C, \Gamma \Sequent
\Delta, !B$, e outra para $!D = !C$, $\Pi = !B, \Lambda$ e $\Lambda
= \Delta$, obtendo !!a{proof}~$\pi_2$ de $!B, !C, \Gamma \Sequent
\Delta, !C$. Eis a !!{proof} de~$!B \land !C, \Gamma \Sequent
\Delta, !B \land !C$:
\[
\AxiomC{}
\RightLabel{$\pi_1$}
\Deduce$!B, !C, \Gamma \fCenter \Delta, !B$
\RightLabel{\LeftR{\land}}
\UnaryInf$!B \land !C, \Gamma \fCenter \Delta, !B$
\AxiomC{}
\RightLabel{$\pi_2$}
\Deduce$!B, !C, \Gamma \fCenter \Delta, !C$
\RightLabel{\LeftR{\land}}
\UnaryInf$!B \land !C, \Gamma \fCenter \Delta, !C$
\RightLabel{\RightR{\land}}
\BinaryInf$!B \land !C, \Gamma \fCenter \Delta, !B \land !C$
\DisplayProof
\]

Para $!A \ident \lforall[x][!B(x)]$, escolha !!a{constant} que não ocorra
em~$\lforall[x][!B(x)]$, $\Gamma$ nem~$\Delta$. Aplique a hipótese de
indução a $!D = !B(c)$, $\Pi = \lforall[x][!B(x)], \Gamma$
e $\Lambda = \Delta$. Isso fornece !!a{proof}~$\pi_1$ de $!B(c),
\lforall[x][!B(x)], \Gamma \Sequent \Delta, !B(c)$, da qual
obtemos:
\[
\AxiomC{}
\RightLabel{$\pi_1$}
\Deduce$!B(c), \lforall[x][!B(x)], \Gamma \fCenter \Delta, !B(c)$
\RightLabel{\LeftR{\lforall}}
\UnaryInf$\lforall[x][!B(x)], \Gamma \fCenter \Delta, !B(c)$
\RightLabel{\RightR{\lforall}}
\UnaryInf$\lforall[x][!B(x)], \Gamma \fCenter \Delta, \lforall[x][!B(x)]$
\DisplayProof
\]
A condição da variável própria em $\RightR{\lforall}$ é satisfeita pela
maneira como escolhemos~$c$.
\end{proof}

\begin{prob}
  Continue a demonstração de \olref{prop:G3c-axioms}, tratando os
  casos em que $!A \ident (!B \lif !C)$ e $!A \ident
  \lexists[x][B(x)]$.
\end{prob}

\end{document}
